\section*{ЗАКЛЮЧЕНИЕ}
\addcontentsline{toc}{section}{ЗАКЛЮЧЕНИЕ}

Разработанная виртуальная файловая система представляет собой полностью функциональный слой абстракции, обеспечивающий прозрачное взаимодействие с данными независимо от их физической природы и местоположения. Унифицированный программный интерфейс, лежащий в основе системы, позволяет выполнять стандартный набор файловых операций над объектами, хранящимися в структурно различных источниках.

Основные результаты работы:

\begin{enumerate}
\item Изучены структуры файловых систем FAT32 и CDFS. Изучена структура блока параметров BIOS(BPB), структура FSInfo. Изучена таблица FAT, структура записи в каталоге. Изучены принципы работы с файлами и каталогами в файловой системе FAT32. Изучен формат данных Little-Endian и Big-Endian. Изучена структура дескрипторов томов, таблицы путей. Изучены принципы работы с файлами и каталогами в файловой системе CDFS.
\item Спроектирована виртуальная файловая система. Спроектированы классы абстрактной файловой системы и абстрактного класса. Спроектированы методы работы с файлом. Спроектированы вспомогательные функции взаимодействия с файловыми системами FAT32 и CDFS.
\item Спроектирован интерфейс виртуальной файловой системы. Спроектированы методы классов файловых систем FAT32 и CDFS.
\item Реализовал виртуальную файловую систему и её интерфейс. Реализованы методы работы с файлом.
\item Реализовал файловые системы FAT32 и CDFS. Реализована инициализация файловых систем. Реализованы методы классов файловых систем.
\item Реализовал командный интерпретатор. Реализованы команды для интерактивного взаимодействия с файловой системой.
\end{enumerate}

Все требования, объявленные в техническом задании, были полностью реализованы, все задачи, поставленные в начале разработки проекта, были также решены.

Готовый рабочий проект представлен в виде библиотек. Библиотеки находится в публичном доступе, поскольку опубликованы в репозитории <<os-pi>>.  
